\documentclass[10pt,letterpaper]{article}
\usepackage[T1]{fontenc}
\usepackage[utf8]{inputenc}
\usepackage[spanish,es-tabla]{babel}
\usepackage{csquotes}
\usepackage{mathpazo}
\usepackage[scaled=0.9]{helvet}
\usepackage{courier}
\usepackage{calc}
\usepackage[top=2cm, bottom=2cm, left=2cm, right=2cm]{geometry}
\usepackage{xcolor}
\usepackage{booktabs}
\usepackage{longtable}
\usepackage{array}
\usepackage{ragged2e}
\usepackage{xurl}
\usepackage[strings]{underscore}
\usepackage{fancyhdr}
\usepackage{sectsty}
\usepackage{tocloft}
\usepackage[colorlinks=true,linkcolor=blue,urlcolor=blue,citecolor=blue,breaklinks=true]{hyperref}
\usepackage{enumitem}

\definecolor{uncpblue}{HTML}{1A237E}
\allsectionsfont{\color{uncpblue}}

\renewcommand{\cftsecleader}{\cftdotfill{\cftsecdotsep}}
\cftpagenumbersoff{section}
\cftpagenumbersoff{subsection}
\cftpagenumbersoff{subsubsection}
\renewcommand{\cftsecfont}{\normalfont}
\renewcommand{\cftsubsecfont}{\normalfont}
\renewcommand{\cftsubsubsecfont}{\normalfont}
\setlength{\parindent}{0pt}
\setlength{\parskip}{4pt}
\setlength{\tabcolsep}{3pt}
\renewcommand{\arraystretch}{0.85}

\pagestyle{fancy}
\fancyhf{}
\fancyhead[L]{\small\textcolor{uncpblue}{\textbf{SGSI-UNCP}}}
\fancyhead[R]{\small\textcolor{uncpblue}{POL-SGSI-06}}
\fancyfoot[C]{\thepage}
\renewcommand{\headrulewidth}{0.4pt}
\renewcommand{\footrulewidth}{0.4pt}
\setlength{\headheight}{14pt}

\setlist{nosep,leftmargin=*,after=\vspace{2pt}}
\setlength{\emergencystretch}{2em}

\newcommand{\makecover}{
\begin{titlepage}
\centering
\vspace*{2cm}
{\Huge\bfseries\color{uncpblue} SGSI-UNCP\par}
\vspace{0.7cm}
{\LARGE Política de Contraseñas y Autenticación Segura\par}
\vspace{0.5cm}
{\large Documento Base del Sistema de Gestión\\de Seguridad de la Información\par}
\vspace{1.5cm}
{\small Universidad Nacional del Centro del Perú\par}
\vfill
\end{titlepage}
}

\begin{document}

\makecover
\setcounter{tocdepth}{2}
\RaggedRight
\tableofcontents
\newpage

\RaggedRight
\section{Política de Contraseñas y Autenticación
Segura}\label{poluxedtica-de-contraseuxf1as-y-autenticaciuxf3n-segura}

\textbf{Organización:} Universidad Nacional del Centro del Perú (UNCP)
\textbf{Aprobado por:} Oficial de Seguridad y Confianza Digital
\textbf{Norma:} ISO/IEC 27001:2022 (Controles A.5.17, A.8.3, A.8.5)
\textbf{Referencia Técnica:} NIST SP 800-63B (Digital Identity
Guidelines) \textbf{Alineamiento:} P-SGSI-03 (Gestión de Accesos),
POL-SGSI-05 (Dispositivos Móviles)

\begin{center}\rule{0.5\linewidth}{0.5pt}\end{center}

\subsection{Objetivo}\label{objetivo}

Establecer los requisitos mínimos de creación, uso, almacenamiento y
rotación de contraseñas en la UNCP, así como las condiciones para el uso
de mecanismos de autenticación multifactor, con el fin de reducir el
riesgo de accesos no autorizados a los sistemas institucionales.

\subsection{Alcance}\label{alcance}

Esta política aplica a todo usuario de los sistemas de información de la
UNCP (personal administrativo, docente, estudiantes, contratistas y
proveedores) que utilice credenciales de acceso para autenticarse en
recursos institucionales.

\begin{center}\rule{0.5\linewidth}{0.5pt}\end{center}

\subsection{Requisitos de
Contraseñas}\label{requisitos-de-contraseuxf1as}

\subsubsection{3.1 Creación y
Complejidad}\label{creaciuxf3n-y-complejidad}

{\setlength{\RaggedRightRightskip}{0pt plus 5em}\small\sloppy \begin{longtable}[]{@{}
  >{\hspace{0pt}\RaggedRight\arraybackslash}p{(\linewidth - 2\tabcolsep) * \real{0.5000}}
  >{\hspace{0pt}\RaggedRight\arraybackslash}p{(\linewidth - 2\tabcolsep) * \real{0.5000}}@{}}
\toprule\noalign{}
\begin{minipage}[b]{\linewidth}\raggedright
Parámetro
\end{minipage} & \begin{minipage}[b]{\linewidth}\raggedright
Requisito Mínimo
\end{minipage} \\
\midrule\noalign{}
\endhead
\bottomrule\noalign{}
\endlastfoot
\textbf{Longitud mínima} & 12 caracteres (se recomienda 14+) \\
\textbf{Longitud máxima} & No debe imponerse límite corto; se aceptan
hasta 128 caracteres \\
\textbf{Composición} & No se exige composición forzada (mayúscula,
minúscula, número, símbolo) según NIST SP 800-63B. Sin embargo, se
recomienda incluir al menos 3 de las 4 categorías \\
\textbf{Caracteres permitidos} & Todos los caracteres ASCII imprimibles,
incluyendo espacio \\
\textbf{Verificación contra listas negras} & La contraseña NO debe estar
en listas de contraseñas comunes comprometidas (Have I Been Pwned,
listas RockYou, etc.) \\
\textbf{Similitud con datos personales} & No debe contener el nombre de
usuario, nombre real, DNI, fecha de nacimiento, o datos institucionales
fácilmente adivinables \\
\end{longtable}\normalsize}

\subsubsection{3.2 Lista Negra de Contraseñas
(Prohibidas)}\label{lista-negra-de-contraseuxf1as-prohibidas}

El sistema de autenticación debe rechazar automáticamente las siguientes
contraseñas (evaluadas por el IdM o política de directorio activo): -
Contraseñas comunes: \texttt{123456}, \texttt{password}, \texttt{admin},
\texttt{admin123}, \texttt{uncp123}, \texttt{universidad},
\texttt{123456789}, etc. - Variantes del nombre de la institución:
\texttt{UNCP2026}, \texttt{univcentro}, \texttt{uncp\_admin}. -
Secuencias de teclado: \texttt{qwerty}, \texttt{asdfgh},
\texttt{zxcvbn}. - Repeticiones: \texttt{aaaaaaaa}, \texttt{12341234}. -
Cualquier contraseña encontrada en violaciones de datos públicas (Have I
Been Pwned).

\subsubsection{3.3 Almacenamiento y
Transmisión}\label{almacenamiento-y-transmisiuxf3n}

\begin{itemize}
\tightlist
\item
  Las contraseñas \textbf{nunca} deben almacenarse en texto plano. Deben
  almacenarse utilizando funciones de hash criptográfico con sal
  (bcrypt, Argon2id, PBKDF2 o similar).
\item
  Las contraseñas deben transmitirse exclusivamente por canales cifrados
  (HTTPS/TLS, VPN).
\item
  Queda prohibido que los sistemas muestren la contraseña en claro en
  ningún momento (pantalla, correo, logs).
\end{itemize}

\subsubsection{3.4 Rotación de
Contraseñas}\label{rotaciuxf3n-de-contraseuxf1as}

\begin{itemize}
\tightlist
\item
  \textbf{No se exige} cambio periódico de contraseña para usuarios
  estándar, salvo que exista sospecha de compromiso (NIST SP 800-63B
  desaconseja el cambio forzado periódico).
\item
  \textbf{Cambio obligatorio} solo en los siguientes casos:

  \begin{itemize}
  \tightlist
  \item
    Sospecha o confirmación de compromiso de la cuenta.
  \item
    Después de la recuperación de una cuenta comprometida.
  \item
    Al reasignar una cuenta a otro usuario.
  \item
    Cambio forzado en la primera autenticación (contraseña temporal).
  \end{itemize}
\item
  Para \textbf{cuentas privilegiadas} (administradores OTI, DBA,
  administradores cloud): rotación cada 90 días.
\end{itemize}

\subsubsection{3.5 Historial}\label{historial}

\begin{itemize}
\tightlist
\item
  No se permite repetir las últimas \textbf{5} contraseñas utilizadas.
\end{itemize}

\begin{center}\rule{0.5\linewidth}{0.5pt}\end{center}

\subsection{Gestión de Credenciales}\label{gestiuxf3n-de-credenciales}

\subsubsection{4.1 Prohibiciones}\label{prohibiciones}

\begin{itemize}
\tightlist
\item
  Las credenciales de acceso son \textbf{personales e intransferibles}.
  Queda prohibido compartir contraseñas con otros usuarios, incluyendo
  compañeros de trabajo o superiores jerárquicos.
\item
  Queda prohibido escribir contraseñas en notas adhesivas visibles,
  debajo del teclado, o en cualquier lugar accesible por terceros.
\item
  Queda prohibido almacenar contraseñas en archivos de texto plano,
  hojas de cálculo, correos electrónicos o documentos sin cifrar.
\end{itemize}

\subsubsection{4.2 Uso de Gestores de
Contraseñas}\label{uso-de-gestores-de-contraseuxf1as}

La UNCP permite y \textbf{recomienda el uso de gestores de contraseñas
empresariales} para el personal administrativo y técnico. El gestor
institucional sugerido debe: - Estar aprobado por la OTI. - Exigir una
contraseña maestra robusta (mínimo 14 caracteres, no reutilizada en
ningún otro servicio). - Tener habilitada la autenticación multifactor
para acceder al cofre de contraseñas. - Permitir el uso compartido
seguro de credenciales entre miembros del equipo (para cuentas de
servicio o compartidas).

\subsubsection{4.3 Cuentas de Servicio y
Compartidas}\label{cuentas-de-servicio-y-compartidas}

\begin{itemize}
\tightlist
\item
  Las cuentas de servicio (no personales) deben tener contraseñas
  generadas aleatoriamente de al menos \textbf{20 caracteres}.
\item
  Las contraseñas de cuentas de servicio deben rotarse cada \textbf{180
  días}.
\item
  Las cuentas compartidas (ej. \texttt{soporte@uncp.edu.pe}) deben estar
  asociadas a un grupo de usuarios nominativos. El acceso a la
  contraseña debe gestionarse a través del gestor de contraseñas
  empresarial.
\end{itemize}

\begin{center}\rule{0.5\linewidth}{0.5pt}\end{center}

\subsection{Autenticación Multifactor
(MFA)}\label{autenticaciuxf3n-multifactor-mfa}

\subsubsection{5.1 Obligatoriedad}\label{obligatoriedad}

{\setlength{\RaggedRightRightskip}{0pt plus 5em}\small\sloppy \begin{longtable}[]{@{}
  >{\hspace{0pt}\RaggedRight\arraybackslash}p{(\linewidth - 4\tabcolsep) * \real{0.3333}}
  >{\hspace{0pt}\RaggedRight\arraybackslash}p{(\linewidth - 4\tabcolsep) * \real{0.3333}}
  >{\hspace{0pt}\RaggedRight\arraybackslash}p{(\linewidth - 4\tabcolsep) * \real{0.3333}}@{}}
\toprule\noalign{}
\begin{minipage}[b]{\linewidth}\raggedright
Tipo de Usuario
\end{minipage} & \begin{minipage}[b]{\linewidth}\raggedright
MFA Exigido
\end{minipage} & \begin{minipage}[b]{\linewidth}\raggedright
Método Principal
\end{minipage} \\
\midrule\noalign{}
\endhead
\bottomrule\noalign{}
\endlastfoot
Administrativo con acceso a sistemas críticos & \textbf{Obligatorio} &
App de autenticación (Microsoft Authenticator, Google Authenticator) \\
Docente (acceso a notas, actas, VPN) & \textbf{Obligatorio} & App de
autenticación \\
Administradores de sistemas (OTI) & \textbf{Obligatorio} & Llavero FIDO2
/ Token físico + app de respaldo \\
Estudiante (acceso a Moodle, correo) & \textbf{Recomendado} & App de
autenticación o SMS \\
Contratista / Proveedor & \textbf{Obligatorio} & App de autenticación \\
\end{longtable}\normalsize}

\subsubsection{5.2 Excepciones al MFA}\label{excepciones-al-mfa}

Solo se permite el acceso sin MFA en los siguientes casos, con
autorización del Oficial de Seguridad: - Cuentas de servicio que no
permiten MFA técnicamente, con controles compensatorios (IP restringida,
certificado de cliente). - Sistemas legacy sin soporte de MFA, con plan
de migración documentado. - Excepción temporal (máximo 30 días) por
pérdida del dispositivo de autenticación.

\begin{center}\rule{0.5\linewidth}{0.5pt}\end{center}

\subsection{Bloqueo de Cuenta}\label{bloqueo-de-cuenta}

{\setlength{\RaggedRightRightskip}{0pt plus 5em}\small\sloppy \begin{longtable}[]{@{}
  >{\hspace{0pt}\RaggedRight\arraybackslash}p{(\linewidth - 2\tabcolsep) * \real{0.5000}}
  >{\hspace{0pt}\RaggedRight\arraybackslash}p{(\linewidth - 2\tabcolsep) * \real{0.5000}}@{}}
\toprule\noalign{}
\begin{minipage}[b]{\linewidth}\raggedright
Parámetro
\end{minipage} & \begin{minipage}[b]{\linewidth}\raggedright
Valor
\end{minipage} \\
\midrule\noalign{}
\endhead
\bottomrule\noalign{}
\endlastfoot
\textbf{Umbral de bloqueo} & 5 intentos fallidos consecutivos \\
\textbf{Duración del bloqueo} & 30 minutos automáticos, o hasta que el
administrador desbloquee manualmente \\
\textbf{Notificación} & Al superar el umbral, se envía una notificación
al usuario y al equipo de seguridad \\
\textbf{Cuentas privilegiadas} & El umbral se reduce a 3 intentos
fallidos. El desbloqueo solo puede realizarse manualmente por el
administrador del sistema \\
\end{longtable}\normalsize}

\begin{center}\rule{0.5\linewidth}{0.5pt}\end{center}

\subsection{Autenticación para Acceso Remoto (Control
A.8.5)}\label{autenticaciuxf3n-para-acceso-remoto-control-a.8.5}

Todo acceso remoto a los sistemas de la UNCP (VPN, escritorio remoto,
consolas de administración cloud) debe cumplir: - MFA obligatorio. -
Sesión con tiempo de expiración por inactividad: máximo \textbf{15
minutos}. - Conexión cifrada (TLS 1.2+, IPsec). - Las cuentas de acceso
remoto deben ser nominativas. No se permiten cuentas genéricas para
acceso remoto.

\begin{center}\rule{0.5\linewidth}{0.5pt}\end{center}

\subsection{Contraseñas Temporales y de Primer
Acceso}\label{contraseuxf1as-temporales-y-de-primer-acceso}

\begin{itemize}
\tightlist
\item
  Las contraseñas temporales (otorgadas en la creación de la cuenta o
  después de un restablecimiento) deben:

  \begin{itemize}
  \tightlist
  \item
    Ser generadas aleatoriamente (mínimo 12 caracteres).
  \item
    Expirar en un máximo de \textbf{24 horas} o en el \textbf{primer
    inicio de sesión exitoso}.
  \item
    Forzar el cambio de contraseña en el primer inicio de sesión.
  \end{itemize}
\item
  No se deben enviar contraseñas temporales por canales no seguros.
  Preferir el envío a través del portal de autoservicio o entrega
  presencial.
\end{itemize}

\begin{center}\rule{0.5\linewidth}{0.5pt}\end{center}

\subsection{Responsabilidades}\label{responsabilidades}

{\setlength{\RaggedRightRightskip}{0pt plus 5em}\small\sloppy \begin{longtable}[]{@{}
  >{\hspace{0pt}\RaggedRight\arraybackslash}p{(\linewidth - 2\tabcolsep) * \real{0.5000}}
  >{\hspace{0pt}\RaggedRight\arraybackslash}p{(\linewidth - 2\tabcolsep) * \real{0.5000}}@{}}
\toprule\noalign{}
\begin{minipage}[b]{\linewidth}\raggedright
Rol
\end{minipage} & \begin{minipage}[b]{\linewidth}\raggedright
Responsabilidad
\end{minipage} \\
\midrule\noalign{}
\endhead
\bottomrule\noalign{}
\endlastfoot
\textbf{Usuario} & Crear contraseñas seguras; no compartir credenciales;
usar MFA; reportar cualquier sospecha de compromiso. \\
\textbf{OTI (Soporte)} & Configurar y mantener las políticas de
contraseñas en el IdM/Directorio Activo; gestionar el restablecimiento
de contraseñas. \\
\textbf{OTI (Seguridad)} & Mantener la lista negra de contraseñas;
auditar la fortaleza de las contraseñas periódicamente; gestionar el
MFA. \\
\textbf{Oficial de Seguridad} & Definir y mantener esta política;
autorizar excepciones; revisar el cumplimiento. \\
\end{longtable}\normalsize}

\begin{center}\rule{0.5\linewidth}{0.5pt}\end{center}

\subsection{Métricas de Cumplimiento}\label{muxe9tricas-de-cumplimiento}

{\setlength{\RaggedRightRightskip}{0pt plus 5em}\small\sloppy \begin{longtable}[]{@{}
  >{\hspace{0pt}\RaggedRight\arraybackslash}p{(\linewidth - 6\tabcolsep) * \real{0.2500}}
  >{\hspace{0pt}\RaggedRight\arraybackslash}p{(\linewidth - 6\tabcolsep) * \real{0.2500}}
  >{\hspace{0pt}\RaggedRight\arraybackslash}p{(\linewidth - 6\tabcolsep) * \real{0.2500}}
  >{\hspace{0pt}\RaggedRight\arraybackslash}p{(\linewidth - 6\tabcolsep) * \real{0.2500}}@{}}
\toprule\noalign{}
\begin{minipage}[b]{\linewidth}\raggedright
Indicador
\end{minipage} & \begin{minipage}[b]{\linewidth}\raggedright
Meta
\end{minipage} & \begin{minipage}[b]{\linewidth}\raggedright
Frecuencia
\end{minipage} & \begin{minipage}[b]{\linewidth}\raggedright
Fuente
\end{minipage} \\
\midrule\noalign{}
\endhead
\bottomrule\noalign{}
\endlastfoot
\% de usuarios con MFA activado (personal) & 100\% & Mensual & IdM /
Directorio Activo \\
\% de cuentas con contraseñas en lista negra & 0\% & Trimestral &
Auditoría de hash de contraseñas \\
Tiempo medio de desbloqueo de cuenta & \textless{} 1 hora & Mensual &
Sistema de tickets \\
Número de intentos de autenticación fallidos & Reporte mensual & Mensual
& SIEM \\
\% de cuentas de servicio con rotación \textless{} 180 días & 100\% &
Trimestral & Sistema PAM \\
\end{longtable}\normalsize}

\begin{center}\rule{0.5\linewidth}{0.5pt}\end{center}

\subsection{Documentos Relacionados}\label{documentos-relacionados}

{\setlength{\RaggedRightRightskip}{0pt plus 5em}\small\sloppy \begin{longtable}[]{@{}
  >{\hspace{0pt}\RaggedRight\arraybackslash}p{(\linewidth - 2\tabcolsep) * \real{0.5000}}
  >{\hspace{0pt}\RaggedRight\arraybackslash}p{(\linewidth - 2\tabcolsep) * \real{0.5000}}@{}}
\toprule\noalign{}
\begin{minipage}[b]{\linewidth}\raggedright
Código
\end{minipage} & \begin{minipage}[b]{\linewidth}\raggedright
Nombre
\end{minipage} \\
\midrule\noalign{}
\endhead
\bottomrule\noalign{}
\endlastfoot
P-SGSI-03 & Gestión de Identidades y Control de Acceso \\
POL-SGSI-02 & Política de Uso Aceptable, Escritorio y Teletrabajo \\
POL-SGSI-05 & Política de Seguridad para Dispositivos Móviles del
Personal \\
D-SGSI-08 & Guía de Implementación de Controles (ISO 27002) \\
\end{longtable}\normalsize}

\begin{center}\rule{0.5\linewidth}{0.5pt}\end{center}

---**

\end{document}
